\section{The boundary exponent}\label{c01-c1-the-boundary-exponent}

We now determine the strip mass needed to normalize chordal walks. The row derivative reduces this problem to a positive polynomial-kernel estimate and a comparison of geometric tails.

\begin{proposition}[Boundary exponent]\label{prop:pv:boundary}
For critical bridges across $h$ layers from a fixed bottom port to all top ports, $b_h=h^{-1/4+o(1)}$. The half-plane arch mass to the right is $1/(2\cos(3\pi/8))$; its height and diameter tails both have exponent $-1/4$.
\end{proposition}

The proof has three stages. A row derivative gives an exact excluded-arch mass. A positive polynomial-kernel estimate gives its power, and a two-direction flux comparison changes the cylinder exclusion into a strip-height tail. This argument uses the row relations and physical eigenline from Lemmas~\ref{lem:pv:row-relations} and~\ref{lem:pv:physical-eigenline}, independently of the polynomial vacuum and separator pairing.

\subsection{Row susceptibility at infinity}\label{c01-c1-row-susceptibility-at-infinity}

Use even widths \(N=2m\). Let \(G_N\) be the half-plane mass, from a fixed bottom port, of arches with positive endpoint displacement \(0<d<N\) whose projections to the cylinder of period \(N\) are simple. The strip flux bound gives \(G_N\le1/(2\cos(3\lambda))\). We will express this deficit through one Taylor coefficient of a row eigenvalue.

For the common eigenline near twist \(\kappa=i\), differentiate its row eigenvalue in \(\eta=\kappa+\kappa^{-1}\) at zero. Call the answer in homogeneous data \(H(\theta)\), with \(z=e^{2iv},\ \theta=2v-3\lambda\); we also write \(H(z)\) when using the transfer argument. It tends to 1 as \(z\to0,\infty\) by triangularity. It is even: at real \(v\) near \(\lambda\), empty-to-empty powers agree with those at \(3\lambda-v\) by reflection of the unoriented loop sum (\(u_1,u_2\) and \(w_1,w_2\) interchange). The simple eigenvalue strictly dominates nearby at both parameters, so they agree there on that eigenvalue and then by continuation. Its denominator may be taken to be \((\cos\lambda+\cos\theta)^N\).

Before making sites homogeneous, transfer fusion gives for the derivative as a function of transfer argument \(z\) both relations
\[
 H(z)+H(Tz)=0,\qquad H(z)+H(t^2z)-H(tz)=0
\]
at every \(z=x_j\) (base eigenvalue 1). Thus the relations have zero jets of orders below \(N\) at the common site 1, by analyticity in sites and coalescence (all the denominators here regular locally; e.g. divide by \(\prod(z-x_j)\) for nearby distinct sites and bound inside a small circle). In \(\theta\) notation the site and its shifts \(t,t^2,T\) correspond to \(-3\lambda,-\lambda,\lambda,3\lambda\). By evenness the jets are
\[
 H^{(2j)}(3\lambda)=0,\qquad
 H^{(2j+1)}(3\lambda)=2H^{(2j+1)}(\lambda)\quad(0\le j<m).
\]

Represent the cosine-rational function as
\[
 H(\theta)=1+\int_{\mathbb R} e^{\theta p}
 \frac{\sinh(\lambda p)U(p^2)+p\cosh(\lambda p)V(p^2)}{\sinh(\pi p)}\,dp
\]
with polynomials of degrees below \(m\). Indeed, near the real angles used here,
\[
 \int_{\mathbb R}\frac{e^{\theta p}\sinh(\alpha p)}{\sinh(\pi p)}\,dp
 =\frac{\sin\alpha}{\cos\theta+\cos\alpha}.
\]
Its successive even \(\theta\)-derivatives and their \(\alpha\)-derivatives give precisely a basis by increasing pole order. The elementary transform follows by taking differences of \({\rm pv}\int e^{u p}/\sinh(\pi p)\,dp=\tan(u/2)\). For real \(|u|<\pi\), take the integral on \(\mathbb R+i\epsilon\), \(0<\epsilon<1\), shift that line by \(i\), and get \(-2i e^{iu}/(1+e^{iu})\) by the pole at \(i\). As \(\epsilon\downarrow0\), the line integral gives the principal value minus \(i\), from \(1/(\pi p)\) at zero.

Put \(x=\lambda p\), \(s=p^2\), and let \(\mu\) be the measure on \(s\ge0\) pushed forward from
\[
 \frac{\cosh(3x)\sinh x}{\sinh(8x)}\,dp.
\]
Define the test and trial spaces
\[
 \begin{aligned}
 \mathcal F_m&=\operatorname{span}\{s^j,s^j p\tanh x:0\le j<m\},\\
 \mathcal G_m&=\operatorname{span}\{s^j,s^j p\coth x:0\le j<m\}.
 \end{aligned}
\]
The even functions of \(p\) here are regarded as functions of \(s\), with their continuous values at zero. The jets say precisely that \(R(s)=U(s)+p\coth x\,V(s)\in\mathcal G_m\) satisfies
\begin{equation}\label{eq:pv:boundary-representer}
 \int_{[0,\infty)}[-R(s)]f(s)\,d\mu(s)=f(0)
 \qquad(f\in\mathcal F_m).
\end{equation}
This uses the identities
\[
 \cosh3x=(2\cosh2x-1)\cosh x,\qquad
 \sinh3x-2\sinh x=(2\cosh2x-1)\sinh x.
\]
At \(z=0\) the first Taylor coefficient is
\begin{equation}\label{eq:pv:boundary-coefficient}
 [z]H=2\sin\lambda\ e^{-3i\lambda}R(-1),
\end{equation}
by shifting the \(p\) contour upwards just past \(i\) at fixed \(\Re\theta=0,\ \Im\theta\to+\infty\): the pole there contributes \(2\sin\lambda\,R(-1)e^{i\theta}\) and the new integral is \(o(|z|)\) by exponential decay.

There is also a path interpretation. Take a homogeneous finite stack of honeycomb rows with one further row of parameter \(v\) above, empty boundaries, and divide by the empty-to-empty value without the further row. This tends to the row eigenvalue uniformly near twist \(i\), by strict spectral dominance and commutation. Its derivative at \(\eta=0\) counts a single essential loop using the further row. This follows by the same arc/turn embedding for angles near \(\lambda\), hence identically. With empty top the loop is either purely horizontal in that row (weight \(o^N\)) or has horizontal intervals connecting bottom ports there, using equal positive numbers of \(u_1,u_2\) and otherwise only factors \(o\) along intervals. Now
\[
 o=\frac{(1-z)(1-t^{-3}z)}{(1-t^2z)(1-Tz)},\qquad
 u_1u_2=o\,\sin^2(2\lambda)/[\sin(2\lambda+v)\sin(3\lambda+v)]\sim-2 e^{5i\lambda}z .
\]
Thus the only other contribution to \([z]\) is a single lower arch between distinct ports plus one horizontal connection in the further row. Such an arch lifts with endpoint displacement strictly below \(N\) in absolute value (a larger displacement interlaces its translate), and of the two horizontal connections exactly one closes it with essential winding. Conversely every arch counted by \(G_N\) contributes, with this choice of starting end unique for each unoriented cylinder arc. The reflected version has the same weight. Passing to the infinite lower stack therefore gives \([z]H=N([z]o-2e^{5i\lambda}G_N)\). This passage on derivative and coefficient is also valid by spectral convergence, uniformly in the twist perturbation and in regular further-row arguments near \(z=0\). Combining \([z]o=2e^{5i\lambda}/(2\cos(3\lambda))\) with \eqref{eq:pv:boundary-coefficient} gives the promised identity:
\begin{equation}\label{eq:pv:excluded-arch-mass}
 \frac1{2\cos(3\lambda)}-G_N=\frac{\sin\lambda}{N}[-R(-1)] .
\end{equation}

\subsection{Positive kernel estimate}\label{c01-c1-positive-kernel-estimate}

By \eqref{eq:pv:excluded-arch-mass}, it remains to prove
\begin{equation}\label{eq:pv:boundary-kernel-power}
 -R(-1)=N^{3/4+o(1)}.
\end{equation}
The two spaces in \eqref{eq:pv:boundary-representer} are different, so that equation does not itself give a positive polynomial kernel. We first express its solution as a positive average of ordinary kernels over two pole sets. A comparison valid for every pole set gives the upper bound. For the lower bound, low poles occur with probability tending to one; on that event a smaller polynomial space admits the reverse comparison.

The partial fractions of \(p\tanh x\) and \(p\coth x\) are respectively \((2/\lambda)s\sum_{a}1/(s+a)\) and \(1/\lambda+(2/\lambda)s\sum_b1/(s+b)\), with \(a=(8l)^2\) at positive half-integers \(l=1/2,3/2,\ldots\), \(b=(8l)^2\) at positive integers. Equivalently the all-pole products \(\prod(1+s/a),\prod(1+s/b)\) equal \(\cosh x,\sinh x/x\). For example integrate \(f(w)/(w^2-x^2)\) for \(f=\tanh,\coth\) on expanding squares spaced away from poles (where \(f\) stays bounded); residues at \(\pm x\) sum to \(f(x)/x\), while the poles of \(f\) each have residue 1 for \(f\). The products follow by integrating logarithmic derivatives from zero.

For subsets \(A,B\) of size \(m\) of the respective pole sets, write \(\pi_A(s)=\prod_{a\in A}(1+s/a)\), similarly \(\pi_B\), and \(dw=d\mu/(\pi_A\pi_B)\). Let \(K_d^w\) denote the ordinary polynomial evaluation kernel through degree \(d-1\): if \(P_j\) are orthonormal for \(w\), then \(K_d^w(u,v)=\sum_{j<d}P_j(u)P_j(v)\).

\begin{lemma}[Positive pole representation]\label{lem:pv:pole-representation}
The weights
\[
 v_A v_B Z_N(w),\qquad v_A=\Delta(A)^2\prod_{a\in A} a^{1-N},\qquad
 Z_N(w)=\frac1{N!}\int\Delta(s_1,\ldots,s_N)^2\prod_i dw(s_i)
\]
with \(\Delta\) Vandermonde have a finite, positive sum over \(A,B\), and hence define a joint probability law after normalization. For this law,
\begin{samepage}
\begin{equation}\label{eq:pv:pole-representation}
 -R(-1)=\mathbb E_{A,B}\frac{K_N^w(0,-1)}{\pi_A(0)\pi_B(-1)}
\end{equation}
and \(\pi_A(0)=1\), \(\sin\lambda/\lambda\le\pi_B(-1)\le1\). In particular all terms in the expectation are positive.
\par
\end{samepage}
\end{lemma}

\begin{proof}
A determinant evaluating the test basis at \(N\) nodes factors as a nonzero signed constant times
\[
 \Delta(s_1,\ldots,s_N)\sum_A\frac{v_A}{\prod_i\pi_A(s_i)};
\]
the trial determinant has the corresponding expansion over \(B\). To obtain these formulas, expand the last \(m\) rows by the partial fractions, discard polynomial-row terms from division, and antisymmetrize the pole choices. One Vandermonde in the poles comes from powers of poles, the other from the Cauchy--Vandermonde determinant with rows \(1,s,\ldots,s^{m-1},1/(s+a)\): clear denominators and evaluate at \(-a\). Truncate pole sums first. After removing the common sign and Vandermonde, the scalar sums increase positively for nodes above \(-16\); the determinant integrals also converge directly by exponential tails.

For column-valued lists \(f,g\) of the test and trial bases, set \(M=\int fg^t\,d\mu\). Andr\'eief's identity~\cite[Eq.~(1.3)]{TracyWidom1998} expresses \(\det M\) as the integral of the product of the two evaluation determinants, divided by \(N!\). Substituting their pole expansions gives a common nonzero constant times \(\sum_{A,B}v_Av_BZ_N(w)\). This proves finiteness, positivity of the normalizing sum, and invertibility of \(M\).

Equation~\eqref{eq:pv:boundary-representer} now gives \(-R(-1)=g(-1)^tM^{-1}f(0)\). Its numerator with denominator \(\det M\) is
\[
 \frac1{(N-1)!}\int
 \det f(s_1,\ldots,s_{N-1},0)
 \det g(s_1,\ldots,s_{N-1},-1)\prod_i d\mu(s_i),
\]
by expanding the final columns and then minors. Inserting the pole expansions again, the same cofactor identity for monomials gives \(Z_N(w)K_N^w(0,-1)\), with the remaining factors \(\pi_A(0)\pi_B(-1)\) in the denominator. The common constants cancel in the ratio, proving \eqref{eq:pv:pole-representation}. These integrals have nonnegative sign after removal of the common constant because the two external nodes lie below the integration nodes. Finally, \(\pi_A(0)=1\), and completing the product over all \(b\)-poles gives
\[
 1\ge\pi_B(-1)\ge\prod_b(1-1/b)=\frac{\sin\lambda}{\lambda}>0.
\]
\end{proof}

\paragraph{Reference weights and uniform comparisons.}
For weights on \(s\ge0\), zeros of the degree-\(d\) orthogonal polynomial are eigenvalues of compression of multiplication by \(s\) to degrees below \(d\): the eigen-equation states orthogonality of \((s-\rho)P\). An increasing likelihood ratio shifts each zero up by min-max, since every Rayleigh quotient increases by monotone tilting. We also compare diagonal kernels by their extremal property,
\[
 K_d^w(u,u)=\sup_{\substack{\deg Q<d\\Q\ne0}}
 \frac{|Q(u)|^2}{\int|Q(s)|^2\,dw(s)}.
\]
Let \(\nu_\alpha\) be the pushforward to \(s=p^2\) of \(dy/\cosh(\pi y)\), where \(y=\alpha p/\pi\) and \(\alpha>0\). Orthonormal polynomials of degree \(j\) in \(s\) have values, up to sign, given by coefficient \(2j\) in
\[
 (1-v^2)^{-1/2}\big((1+v)/(1-v)\big)^{iy}.
\]
Indeed integrating a product of two generating functions gives \(1/(1+uv)\) by \(\int e^{iay}dy/\cosh(\pi y)=1/\cosh(a/2)\) (shift contour by \(i\)); use just even indices. At \(s=0\) the magnitudes are of order \((j+1)^{-1/2}\); at \(s=-1\) and \(r=\alpha/\pi\in(1/2,1)\), of order \((j+1)^{r-1/2}\). Indeed the coefficients needed there are those of \((1-v)^{-1/2-r}(1+v)^{r-1/2}\). The first factor has increasing coefficients of order \((k+1)^{r-1/2}\) with ratio tending to 1 (binomial series); the second has summable coefficients with nonzero sum, giving the claim by convolution.

The all-poles lower bound for \(w\) comes from density \(w_\infty(p)=x(2\cosh2x-1)/\sinh(8x)\) before pushing forward. Its ratio to the reference density at \(\alpha=6\lambda\) is increasing on \(p\ge0\): up to a constant this ratio is
\(x(2-1/\cosh2x)/(\tanh2x+\tanh6x)\). Removing poles only increases the upward tilt. Conversely \(w\) relative to reference \(\alpha=2\lambda\) has decreasing ratio since \(\mu\) times \(\cosh2x\) has density proportional to \((2\cosh2x-1)/\cosh4x\), decreasing. It follows that
\[
 K_N^w(0,-1)\le C\sqrt{\log(2N)}\ N^{3/4}.
\]
(Indeed \(w\) dominates a constant times the \(6\lambda\) reference uniformly, so the two diagonal kernel values by the supremum characterization are at most constant times the reference sums of squared coefficients; now use Cauchy--Schwarz.)
Moreover, for every orthogonal polynomial \(P_j\) of degree \(j\) for any of the conditional weights \(w\),
\[
 1\le P_j(-1)/P_j(0)\le C(j+1)^{3/4}
\]
by positivity and comparison of zeros with \(\alpha=6\lambda\).

The upper kernel bound holds for every pole set. The lower bound will use pole sets that contain all low poles; the next lemma shows that these account for asymptotically all of the probability in \eqref{eq:pv:pole-representation}.

\begin{lemma}[Low-pole saturation]\label{lem:pv:pole-saturation}
Fix \(0<\delta<1/2\) and put \(J=N^{1-\delta}\). Under the law of Lemma~\ref{lem:pv:pole-representation}, let \(\mathcal E_J\) be the event that \(A\) and \(B\) each contain every pole in their respective sets with index \(l\le J\). Then
\[
 \mathbb P(\mathcal E_J)=1-o(1).
\]
For either subset, the probability of missing such a pole is \(o(1)\) uniformly conditional on the other subset.
\end{lemma}

\begin{proof}
Fix the other subset. If an index \(j\le J\) is missing, replace the smallest occupied \(k\ge2J\) in the chosen subset by \(j\). Such a \(k\) exists for all large \(N\), since the subset has \(N/2\) elements and only \(O(J)=o(N)\) indices lie below \(2J\). For fixed \(k\) the replacement is injective. The ratio of old to new weight is
\[
 (k/j)^2\prod_{l\ne k}\left(\frac{k^2-l^2}{j^2-l^2}\right)^2
 \ \mathbb E\prod_{i=1}^N\frac{s_i+(8j)^2}{s_i+(8k)^2}
\]
(the product is over other indices of that subset, and the expectation is under the new polynomial ensemble proportional to \(\Delta^2\prod_i dw_{\rm new}(s_i)\)). The first factors cost at most \(\exp(CJ)(k/J)^{CJ}\). Larger indices than \(k\) cost nothing; by minimality of \(k\), smaller occupied indices lie below \(2J\), where
\[
 \sum_l\big[\log^+(J/l)+\log^+(J/|j-l|)\big]\le CJ
\]
by unit spacing, excluding \(j\).

The expectation is the ratio of degree-\(N\) monic orthogonal polynomial values at \(-(8j)^2,-(8k)^2\) for weight \(w_{\rm new}/(s+(8k)^2)\). Indeed, averaging \(\prod_{i=1}^N(u-s_i)\) against a normalized \(\Delta^2\)-ensemble gives its monic orthogonal polynomial. Its integral against \(u^r\), \(r<N\), for the one-point base weight vanishes: it contains \(\Delta(s_1,\ldots,s_N,u)\), while antisymmetrizing the other factor \(\Delta(s_1,\ldots,s_N)u^r\) gives zero.

The extra denominator \(s+(8k)^2\) preserves the decreasing likelihood ratio relative to \(\nu_{2\lambda}\). The zeros are therefore no larger than those for this reference weight. Since \((s+(8j)^2)/(s+(8k)^2)\) increases in \(s\), zero comparison bounds the polynomial ratio above by the reference ratio. Define
\[
 c_M(R)=[v^M](1-v)^{-1/2-R}(1+v)^{-1/2+R},\qquad M=2N.
\]
We have bounded the old-to-new weight ratio, uniformly in both subsets apart from \(j,k\), by
\begin{equation}\label{eq:pv:pole-replacement-bound}
 \exp(CJ)(k/J)^{CJ}\frac{c_M(2j)}{c_M(2k)}.
\end{equation}

We next show that \eqref{eq:pv:pole-replacement-bound} sums to \(o(1/J)\) over \(k\ge2J\), uniformly in \(j\le J\). Coefficients increase positively with \(R\), since the generating function contains \(\exp(2R\,\mathrm{atanh}\,v)\). Replacing \(2k\) by \(2k-1/2\) gives the lower bound
\[
 c_M(2k)\ge\binom{M+2k-1}{2k-1}.
\]
For \(2J\le k\le N^{1-\delta/2}\), bounding the numerator on radius \(1-k/N\) gives \(\exp(O(k))(CN/k)^{2j+1}\), while the denominator is at least \((N/k)^{2k-1}\). The logarithms of the replacement cost are \(O(k)\), whereas
\[
 (2k-1-2j-1)\log(N/k)\gtrsim\delta k\log N,
\]
which absorbs that cost and makes the sum negligible.

For \(N^{1-\delta/2}<k<N\), the denominator already grows as \(\exp(c_\delta N^{1-\delta/2}\log N)\), while costs and numerator, now estimated on radius \(1-J/N\), are at most \(\exp(CJ\log N)\). Finally, for \(k\ge N\), the binomial gives \(\exp(cN)(k/N)^M\), sufficient also on summing to infinity. Injectivity for fixed \(k\) bounds the conditional probability of missing each \(j\) by \(o(1/J)\). The union bound over \(O(J)\) low indices, and then over the two subsets, proves the lemma without any independence assumption.
\end{proof}

\paragraph{From saturation to the lower kernel bound.}
On \(\mathcal E_J\), the weight ratio to \(w_\infty\) is bounded by \(\exp(Cp^2/J)\), since only poles with indices above \(J\) can be missing. Fix \(\alpha'<6\lambda\) close to \(6\lambda\). On \(|p|\le\epsilon J\), choosing \(\epsilon>0\) small makes \(w\) bounded above by a constant times \(\nu_{\alpha'}\). This pointwise comparison need not hold on the whole line. Put \(d=\lfloor N^{1-2\delta}\rfloor\). What the kernel estimate requires, and what does hold uniformly for all saturated pole sets, is the bound
\begin{equation}\label{eq:pv:saturated-quadratic-form}
 \int |Q(s)|^2\,dw(s)
 \le C\int |Q(s)|^2\,d\nu_{\alpha'}(s)
 \qquad(\deg Q<d).
\end{equation}

To include the tails, expand \(Q\) in the reference orthonormal polynomials. Their diagonal kernel satisfies
\[
 |Q(p^2)|^2
 \le \|Q\|_{L^2(\nu_{\alpha'})}^2
       \exp(C_r d+Cr|p|),
\]
using the generating function on a small fixed radius \(r\). The density of \(w\) before pushing forward is always at most \(Ce^{-4\lambda|p|}\). Choose \(r\) small enough to leave exponential decay in \(|p|\); the integral over \(|p|>\epsilon J\) is then at most
\[
 C\exp(C_r d-c\epsilon J)\,\|Q\|_{L^2(\nu_{\alpha'})}^2.
\]
Since \(d=o(J)\), this proves \eqref{eq:pv:saturated-quadratic-form} together with the comparison on the central interval.

The zeros of every \(P_j\) lie in \((0,\infty)\), so \(P_j(0)P_j(-1)>0\). The bound \(P_j(-1)/P_j(0)\le C(j+1)^{3/4}\) and the extremal kernel comparison from \eqref{eq:pv:saturated-quadratic-form} therefore give, on \(\mathcal E_J\),
\[
 K_N^w(0,-1)\ge C^{-1}d^{-3/4} K_d^w(-1,-1)
 \ge c\,d^{2\alpha'/\pi-3/4}.
\]
Here the second inequality uses the reference sum of squared values at \(-1\). The event has probability \(1-o(1)\) by Lemma~\ref{lem:pv:pole-saturation}, and the denominator in \eqref{eq:pv:pole-representation} is bounded above and below. Together with the uniform upper bound, letting \(\delta\downarrow0\) and \(\alpha'\uparrow6\lambda\) proves \eqref{eq:pv:boundary-kernel-power}.

\subsection{Comparing the strip tails}\label{c01-c1-comparing-the-strip-tails}

The total half-plane arch weight to the right is now exactly \(1/(2\cos(3\lambda))\): this was the upper bound by flux, and \(G_N\) tends to it. Equations~\eqref{eq:pv:excluded-arch-mass} and~\eqref{eq:pv:boundary-kernel-power} give
\[
 d_N:=\frac1{2\cos(3\lambda)}-G_N=N^{-1/4+o(1)}
\]
for even \(N\). This is the mass excluded by displacement or failure of simple cylinder projection. We first transfer its exponent to a spatial maximum, and then use a second tilted direction to force the same exponent for height.

\paragraph{Two tilted flux identities.}

Use these positive-endpoint arches above a horizontal bottom from coordinate 0, and write their endpoint as \(D>0\), height maximum as \(Y\), and
\[
 M_+=\max(A x+B y),\quad M_-=\max(A x-B y),\qquad A=\sqrt3/2,\ B=1/2
\]
in planar coordinates. Write \(P\) for their unnormalized weighted law. At large positive levels \(r\) on the respective tiling grids,
\[
 P(M_\pm-A D\ge r)=q_\pm P(M_\pm\ge r),\quad
 q_\pm=\frac{\sin(\sigma(\pi-\phi_\pm))}{\sin(\sigma(\pi+\phi_\pm))},\qquad
 (\phi_+,\phi_-)=(\pi/3,2\pi/3).
\]
Indeed, in a large convex exhaustion add the cut \(Ax\pm By\le r\) to the domain. Subtracting the two identities of Lemma~\ref{lem:pv:flux} puts the phased sum over old exits crossing that cut on the ray \(e^{-i\sigma\phi_\pm}\) by convexity. Tiling cuts of differing directions never share ports. Old non-bottom exits vanish in weight in the exhaustion by the now saturated arch mass: bottom real parts increase to 1 in total, and all real parts control absolute weights. Bottom exits crossing the cut have phases \(e^{-i\sigma\pi}\) to the right and \(e^{i\sigma\pi}\) to the left, with the latter events corresponding by path reversal and translation to the shifted maxima above. Taking the imaginary part after rotating the ray to angle zero gives the identity. Note \(0<q_-<q_+<1\).

\paragraph{From cylinder exclusion to spatial extent.}
Put \(U(r)=P(M_+\ge r)\). First, \(P(D\ge r)\le r^{-1/4+o(1)}\) by \(d_N\). Conversely, an arch excluded at period \(N\) has horizontal range at least of order \(N\): either its displacement is at least \(N\), or two vertices that differ by a nonzero multiple of \(N\) prevent simple projection. Its maximum horizontal coordinate or the negative of its minimum is therefore at least of order \(N\). Reflection with reversal, replacing \(x\) by \(D-x\), turns the far-left case into the far-right case and preserves the positive-endpoint arch law. Thus \(U(r)\ge r^{-1/4-o(1)}\). Both comparisons use even periods on the appropriate side by constant factors.

Splitting on large endpoint and using the plus ray identity gives, for any fixed \(\rho<1\) close to 1,
\[
 U(r)\le P(D\ge c_\rho r)+q_+ U(\rho r)
\]
for large \(r\), rounding lines on grids of bounded spacing. Choose \(\rho\) sufficiently close to 1 that \(q_+\rho^{-1/4}<1\); iteration gives \(U(r)=r^{-1/4+o(1)}\). In particular the weight of arches reaching distance \(r\) is at most \(r^{-1/4+o(1)}\), by covering both sides using also the reflected maxima.

\paragraph{From spatial extent to height.}
To get a lower height tail, use \(M_-\le M_+\le M_-+2BY\) and both ray identities. Again for any \(\rho<1\) close to 1,
\[
 U(r)\le q_+^{-1} P(Y\ge c'_\rho r)+(q_-/q_+) U(\rho r).
\]
Indeed start at a plus-grid level just below \(r\), use its shifted-max identity, and except when \(Y\) is large compare this event to the analogous shifted minus event at a grid level just above \(\rho r\). Choose \((q_-/q_+)\rho^{-1/4}<1\) and iterate for \(k\) steps with \(\rho^k r\) of order \(r^{1-\epsilon}\). The remainder \((q_-/q_+)^kU(\rho^k r)\) is negligible compared to \(U(r)\) by the latter's exponent. The sum of the height-tail terms is at most a constant times \(P(Y\ge c r^{1-\epsilon})\), by monotonicity and the geometric coefficients. Hence
\[
 P(Y\ge c r^{1-\epsilon})\ge r^{-1/4-o(1)}.
\]
As \(\epsilon\) is arbitrary, this gives the matching height exponent, and therefore the diameter exponent as well. Finally \eqref{eq:pv:strip-flux} gives \(b_h=1-\cos(3\lambda)A_h\). The missing arch mass is the height tail, up to the constant and rounding of heights, proving Proposition~\ref{prop:pv:boundary}.
